%% SCRAP: experiments/02-experiments/heartbeat-doe/corrected-design
%% SOURCE: docs/working/experiments/02-experiments/heartbeat-doe/corrected-design.md
%% STATUS: WORKING
%% FITS: experiments/ch-heartbeat
%% EDITORIAL: lifted — prose rewritten to press voice

\section{Heartbeat DoE — Corrected Design}
\label{sec:heartbeat-doe-corrected-design}

\subsection{Critical Discovery}

The initial Stage~2 design measured \emph{jitter in a fixed heartbeat
interval}, implicitly assuming the tick rate was constant at approximately
1~ms. This assumption is incorrect. The heartbeat tick interval
\texttt{tick\_ns} is a dynamic variable that the inference engine adjusts in
response to workload intensity:

\begin{lstlisting}[language=C]
uint64_t tick_ns;  /* Nanoseconds between ticks -- VARIES with load */
\end{lstlisting}

Stage~1 was conducted with \texttt{HEARTBEAT\_THREAD\_ENABLED=0}, which fixed
the interval and disabled the adaptive mechanism entirely. Stage~2 must
re-enable the heartbeat thread and measure the quality of its adaptive
modulation.

\subsection{The Adaptive Feedback Loop}

The heartbeat acts as a pacing mechanism, not merely a timer. Under high
workload, the inference engine increases \texttt{tick\_ns} — slowing the
heartbeat — to allocate more CPU time for physics-parameter tuning. The
expected signal chain is:

\begin{enumerate}
  \item Inference engine detects increased workload intensity.
  \item \texttt{tick\_ns} is increased (heartrate decreases).
  \item The physics engine receives more tuning time per tick.
  \item The system adapts and eventually converges to a new steady state.
  \item \texttt{tick\_ns} stabilises at the new optimal rate.
\end{enumerate}

\subsection{Corrected Primary Metric}

The primary discriminating metric for Stage~2 is the Pearson correlation
between per-tick workload intensity and per-tick heartrate:

\begin{equation}
  \rho = \frac{\operatorname{Cov}(W, H)}{\sqrt{\operatorname{Var}(W)\operatorname{Var}(H)}}
\end{equation}

where $W_t$ is the dictionary-lookup count in tick~$t$ and $H_t = 10^9 /
\texttt{tick\_ns}_t$ is the heartrate in Hz. A strong positive $\rho$
indicates that the configuration's heartbeat responds proportionally to load.
Target: $\rho > 0.75$.

\subsection{Metric Summary}

\begin{center}
\begin{tabular}{lll}
\toprule
Metric & Purpose & Target \\
\midrule
\texttt{load\_to\_heartrate\_correlation} & Load--heartrate coupling & $> 0.75$ \\
\texttt{heartrate\_responsiveness\_score} & Composite adaptation score & $> 80/100$ \\
\texttt{heartrate\_convergence\_time\_ms} & Time to stable rate & $< 5{,}000$ ms \\
\texttt{mean\_heartbeat\_rate\_hz} & Baseline operating rate & $\approx 1{,}000$ Hz \\
\texttt{heartbeat\_rate\_cv} & Rate variation & $0.05$--$0.20$ \\
\bottomrule
\end{tabular}
\end{center}

\subsection{Implementation Requirements}

Four code changes are required:

\begin{enumerate}
  \item Add \texttt{HeartbeatRateSample} struct to \texttt{include/vm.h}:
        fields for tick number, \texttt{tick\_ns}, workload operations this
        tick, wall-clock timestamp, and physics state (\texttt{hot\_word\_count},
        \texttt{total\_heat}).
  \item Extend the \texttt{VM} struct with a pointer to a rate-sample ring
        buffer and a sample count.
  \item Call \texttt{capture\_heartbeat\_rate\_sample()} in
        \texttt{heartbeat\_thread\_main()} before each \texttt{nanosleep()}.
  \item Extend \texttt{DoeMetrics} and \texttt{metrics\_from\_vm()} to compute
        correlation, responsiveness score, and convergence time from the
        collected samples.
\end{enumerate}

\subsection{Expected Discrimination}

The corrected design is expected to reveal clear separation among the five
elite configurations on the load-correlation axis. A configuration with
$\rho < 0.2$ fails to adapt its heartrate to load and should not be selected
as the golden configuration regardless of its static performance score.

The golden configuration — the one used as the MamaForth production baseline
— is identified as the configuration with the highest composite heartrate
responsiveness score, subject to a minimum correlation threshold of
$\rho \geq 0.75$.

%% TODO(bob): confirm whether corrected-design experiment was ever executed;
%%   update status to HISTORICAL once results are in hand

